\newcommand{\ExplainArray}{\textbf{Array} är en förändringsbar, indexerbar samling med fast storlek där alla element är av samma typ och nås med nollbaserat index inom hakparenteser; Array är en grundläggande datastruktur med O(1)-åtkomst men fast storlek som inte kan ändras efter skapande.}
\newcommand{\ExplainRange}{\textbf{Range} är en oföränderlig, lat samling i Scala som representerar ett intervall av heltal med start, slut och steglängd; den lagrar inte alla värden utan beräknar dem vid behov, vilket gör den mycket minneseffektiv för stora intervall.}
\newcommand{\ExplainVector}{\textbf{Vector} är Scalas standardval för oföränderliga, indexerbara sekvenssamlingar; den har effektiv (nästan O(1)) åtkomst, uppdatering och tillägg tack vare en intern trädbaserad (trie) struktur och rekommenderas som default-sekvenssamling i Scala.}
\newcommand{\ExplainAbstract}{Att \textbf{abstrahera} innebär att man identifierar gemensamma egenskaper och skapar nya begrepp som döljer detaljer och förenklar kodningen; abstraktioner gör kod mer återanvändbar och lättare att resonera om.}
\newcommand{\ExplainAbstractClass}{En \textbf{abstrakt klass} kan ha konstruktorparametrar och konkreta medlemmar, men kan inte instansieras direkt och kan inte mixas in; den används som basklass i en arvshierarki där subklasser implementerar de abstrakta medlemmarna.}
\newcommand{\ExplainAbstractMember}{En \textbf{abstrakt medlem} deklareras utan implementation i en klass eller trait och måste implementeras av konkreta subtyper innan de kan instansieras; abstrakta medlemmar specificerar ett kontrakt som subtyper är skyldiga att uppfylla.}
\newcommand{\ExplainAlgorithm}{En \textbf{algoritm} är en ändlig, väldefinierad sekvens av steg som löser ett givet problem och terminerar med ett korrekt resultat; en algoritm kan implementeras i kod oberoende av programmeringsspråk och analyseras med avseende på korrekthet, tids- och minneskomplexitet.}
\newcommand{\ExplainAnonymousClass}{En \textbf{anonym klass} saknar namn och skapas direkt vid instansiering med 'new' följt av ett trait- eller klassnamn och ett block med extra implementation.}
\newcommand{\ExplainAnonymousFunction}{En \textbf{anonym funktion}, även kallad lambda, är en funktion utan namn som definieras inline med syntaxen '(parametrar) => uttryck'; den kan tilldelas en variabel, skickas som argument till en högre ordningens funktion, eller returneras som värde.}
\newcommand{\ExplainAssignment}{En \textbf{tilldelning}, skriven med '=', ändrar värdet på en förändringsbar variabel (var) till ett nytt värde; tilldelning är inte möjlig på en oföränderlig variabel (val) efter att den initialiserats, vilket uppmuntrar till funktionell programmeringsstil.}
\newcommand{\ExplainAttribute}{Ett \textbf{attribut} är en variabel definierad som del av en klass eller ett objekt och representerar en del av dess tillstånd; attribut kan vara oföränderliga (val) eller föränderliga (var), och görs ofta privata för att skydda mot godtyckliga förändringar}
\newcommand{\ExplainBaseType}{En \textbf{bastyp}, även kallad superklass eller supertrait, är den mest generella typen i en arvshierarki och definierar de egenskaper och metoder som delas av alla subtyper; variabler av bastyp kan referera till instanser av vilken subtyp som helst.}
\newcommand{\ExplainBlock}{Ett \textbf{block} är en sekvens av satser och uttryck (i Scala 3 avgränsade av indentation) som kan innehålla lokala namnbindningar synliga endast inuti blocket; värdet av ett block är värdet av dess sista uttryck.}
\newcommand{\ExplainBoolean}{En \textbf{boolesk} typ kan bara anta värdena 'true' (sann) eller 'false' (falsk) och används i villkorliga uttryck och logiska operationer med operatorerna && (och), || (eller) och ! (negation); booleska uttryck styr programmets kontrollflöde.}
\newcommand{\ExplainCallByName}{Vid \textbf{namnanrop} skickas ett argument som ett fördröjt uttryck som evalueras varje gång det används inuti funktionen; detta möjliggör skapande av egna kontrollstrukturer, och deklareras med '=>' före parametertypen.}
\newcommand{\ExplainCallByValue}{Vid \textbf{värdeanrop} evalueras argumentet exakt en gång innan funktionen anropas och det beräknade värdet binds till parametern; detta är standardbeteendet i Scala och de flesta andra programmeringsspråk.}
\newcommand{\ExplainCaseClass}{En \textbf{case-klass} genererar automatiskt ett kompanjonsobjekt med apply-metod (slipper 'new'), equals och hashCode baserade på fältinnehåll, en läsbar toString, en copy-metod för att skapa modifierade kopior, samt stöd för mönstermatchning med unapply.}
\newcommand{\ExplainClass}{En \textbf{klass} är en mall som beskriver struktur och beteende för objekt av en viss typ; den kan ha konstruktorparametrar, attribut och metoder, och varje instans skapad med 'new' får sitt eget tillståndsminne men delar metodimplementationerna.}
\newcommand{\ExplainClassParameter}{En \textbf{klassparameter} deklareras i klassens rubrik och binds till argumentet som ges vid konstruktion; om den föregås av 'val' eller 'var' blir den automatiskt ett synligt attribut, annars är den ett privat konstruktorargument.}
\newcommand{\ExplainCollection}{En \textbf{samling} är en datastruktur som innehåller ett antal element av samma typ och erbjuder operationer för att lägga till, ta bort, söka och transformera element; Scala erbjuder ett rikt samlingsbibliotek med både föränderliga och oföränderliga varianter.}
\newcommand{\ExplainCollectionLibrary}{Scalas \textbf{samlingsbibliotek} innehåller ett rikt urval av samlingstyper med ett enhetligt gränssnitt uppdelat i föränderliga (scala.collection.mutable) och oföränderliga (scala.collection.immutable) varianter; de vanligaste är Vector, List, Set, Map och ArrayBuffer.}
\newcommand{\ExplainColumn}{\textbf{Kolonn} är ett alternativt ord för kolumn och används i samband med matriser och tabeller för att beteckna en vertikal sekvens av värden som delar samma kolumnindex.}
\newcommand{\ExplainColumnVector}{En \textbf{kolumnvektor} är en matris med dimensionen m×1 som innehåller m element i en enda vertikal kolumn; den används inom linjär algebra och numerisk beräkning och är ett specialfall av en matris.}
\newcommand{\ExplainCompanionObject}{Ett \textbf{kompanjonsobjekt} har samma namn som en klass, definieras i samma fil, och har ömsesidig tillgång till klassens privata medlemmar; det används ofta för fabriksmetoder, konstanter och hjälpfunktioner kopplade till klassen.}
\newcommand{\ExplainCompile}{Att \textbf{kompilera} innebär att en kompilator analyserar källkoden, utför typkontroll och andra statiska analyser, och översätter den till en körbar form såsom bytekod för JVM; eventuella fel rapporteras som kompileringsfel.}
\newcommand{\ExplainCompileError}{Ett \textbf{kompileringsfel} uppstår när kompilatorn inte kan tolka eller typkontrollera koden och avbryter kompileringen med ett felmeddelande som anger plats och orsak; kompileringsfel måste åtgärdas innan programmet kan köras.}
\newcommand{\ExplainCompiler}{En \textbf{kompilator} är ett program som läser källkod, analyserar den för fel och översätter den till maskinkod som datorn kan köra; kompilatorn rapporterar kompileringsfel om koden inte är korrekt.}
\newcommand{\ExplainComputer}{En \textbf{dator} är en elektronisk enhet som behandlar data enligt lagrade instruktioner, och utför beräkningar, minnesmanipulationer och logiska operationer för att lösa problem eller utföra uppgifter. En dator består av ev en centralprocessorenhet (CPU), in- och utdataenheter, samt ett minne som lagrar instruktioner och data representerade som heltal.}
\newcommand{\ExplainConstructor}{En \textbf{konstruktor} är den kod som körs när en ny instans skapas med 'new'; i Scala utgörs primärkonstruktorn av klassens rubrik och kropp, och den allokerar minne för instansens tillståndsvariabler och exekverar initialiseringskoden.}
\newcommand{\ExplainControlStructure}{En \textbf{kontrollstruktur} kan ge alternativa vägar genom koden. Exempel på kontrollstrukturer är if-uttryck för logiska val, for-uttryck för att gå igenom sekvenser, och while-satser för att repetera kod så länge ett villkor är uppfyllt.}
\newcommand{\ExplainCounting}{\textbf{Registrering} är ett algoritmiskt mönster som räknar hur många element i en samling som uppfyller ett visst kriterium; det implementeras t.ex. med samlingsmetoden 'count' eller med en räknarvariabel som ökas i en loop.}
\newcommand{\ExplainDataStructure}{En \textbf{datastruktur} är en organiserad samling av data där elementen kan lagras, nås och modifieras på ett väldefinierat sätt; valet av datastruktur, t.ex. array, länkad lista eller träd, påverkar programmets prestanda avsevärt.}
\newcommand{\ExplainDefaultArgument}{Ett \textbf{defaultargument} är ett förvalt värde för en parameter som används automatiskt om inget argument anges vid anropet; detta minskar behovet av överlagring, gör API:er mer flexibla och förbättrar läsbarheten.}
\newcommand{\ExplainDeserialize}{Att deserialisera innebär att man läser en teckensekvens eller byteström (t.ex. JSON, XML eller binärdata) och rekonstruerar de ursprungliga objekten i minnet med rätt typer och tillstånd; deserialisering är inversen till serialisering.}
\newcommand{\ExplainDotNotation}{\textbf{Punktnotation} används för att komma åt ett objekts synliga medlemmar genom att skriva objektreferensen, en punkt, och sedan medlemmens namn, t.ex. 'obj.metod()' eller 'obj.fält'; metoder kan anropas utan parenteser om de saknar parametrar.}
\newcommand{\ExplainDynamicBinding}{\textbf{Dynamisk bindning} innebär att det vid ett metodanrop är objektets körtidstyp, inte den statiska deklarerade typen, som avgör vilken implementering av metoden som faktiskt exekveras; detta är grunden för polymorfism och möjliggörs av överskuggning.}
\newcommand{\ExplainElement}{Ett \textbf{element} är ett enskilt objekt som ingår i en datastruktur eller samling; element kan nås via index i sekvenser, via nyckel i mappningar, eller via iteration med t.ex. en for-loop eller en högre ordningens funktion.}
\newcommand{\ExplainEnumeration}{En \textbf{enumeration} är en datastruktur som består av ett begränsat antal namngivna värden i bestämd ordning; den används för att representera en uppräknelig mängd alternativ med tydliga namn och kan förbättra typkontroll och läsbarhet.}
\newcommand{\ExplainExport}{\textbf{Export}-satsen i Scala 3 gör att utvalda medlemmar från ett inre objekt eller en delegation blir synliga som direkta medlemmar i det omgivande objektet, vilket förenklar gränssnitt och möjliggör transparent delegering.}
\newcommand{\ExplainExpression}{Ett \textbf{uttryck} är en kombination av värden, operatorer och funktionsanrop som evalueras och producerar ett resultat med ett specifikt värde och en specifik typ; i Scala är nästan allt uttryck, inklusive if-satser, match och block.}
\newcommand{\ExplainFactoryMethod}{En \textbf{fabriksmetod} är en funktion, ofta placerad i ett kompanjonsobjekt, som skapar och returnerar instanser av en klass utan att anroparen behöver använda 'new' direkt; den kan validera argument, välja en lämplig subtyp eller returnera en cachad instans.}
\newcommand{\ExplainFloatingpointNumber}{Ett \textbf{flyttal} representerar ett reellt tal med begränsad precision enligt IEEE 754-standarden; Scala har typerna Float (32 bitar) och Double (64 bitar), och avrundningsfel kan uppstå i beräkningar vilket gör att exakt jämförelse med == bör undvikas.}
\newcommand{\ExplainForStatement}{En \textbf{for-sats} itererar över ett intervall eller en samling och exekverar kroppen ett bestämt antal gånger; i Scala kan 'for' användas som sats (utan 'yield') för sidoeffekter, eller som uttryck med 'yield' för att producera en ny samling.}
\newcommand{\ExplainFunction}{En \textbf{funktion} tar noll eller fler inparametrar och beräknar ett returvärde; i Scala är funktioner förstklassiga värden som kan tilldelas variabler, skickas som argument och returneras från andra funktioner, vilket stöder funktionell programmeringsstil.}
\newcommand{\ExplainFunctionBody}{\textbf{Funktionskropp}en är den kod som exekveras när funktionen anropas; den kan vara ett enstaka uttryck efter '=' eller ett block som innehåller lokala variabler och satser vars sista rad ger returvärdet.}
\newcommand{\ExplainFunctionHeader}{\textbf{Funktionshuvud}et börjar med nyckelordet 'def' följt av funktionsnamn, eventuella typparametrar inom hakparenteser, en eller flera parameterlistor med namn och typer, samt en valfri returtypsannotering efter kolon.}
\newcommand{\ExplainGeneric}{En \textbf{generisk} klass, trait eller funktion har en eller fler abstrakta typparametrar skrivna inom hakparenteser som konkretiseras med faktiska typer vid användning; detta möjliggör återanvändbar kod som fungerar korrekt för många olika typer, t.ex. 'Vector[Int]' eller 'Vector[String]'.}
\newcommand{\ExplainGetter}{En \textbf{getter} är en metod som ger indirekt läsåtkomst till ett privat attribut; i Scala definieras den typiskt som en 'def' eller 'val' med samma namn som attributet, utan 'private'-modifieraren, och följer principen om enhetlig access.}
\newcommand{\ExplainImplementation}{En \textbf{implementation} är en konkret realisering av en algoritm, ett gränssnitt eller en abstrakt specifikation i ett specifikt programmeringsspråk; samma algoritm kan ha flera implementationer med olika tids- och minneskomplexitet.}
\newcommand{\ExplainImport}{En \textbf{import}-sats gör ett namn tillgängligt i det lokala omfånget utan att hela den kvalificerade sökvägen behöver skrivas varje gång; i Scala 3 kan import placeras var som helst i koden och man kan importera flera namn eller använda wildcard med '_'.}
\newcommand{\ExplainInheritance}{Med \textbf{arv} kan en typ (subtyp) ärvas från en annan typ (supertyp), vilket innebär att subtypen får tillgång till supertypens medlemmar och kan användas där supertypen förväntas; det uttrycker relationen 'är en' och stödjer polymorfism och specialisering.}
\newcommand{\ExplainInstance}{En \textbf{instans} är ett konkret objekt skapat från en klass med hjälp av 'new' (eller en fabriksmetod); varje instans har sitt eget minne för tillståndsvariabler men delar metodimplementationerna med alla andra instanser av samma klass.}
\newcommand{\ExplainKey}{En \textbf{nyckel} är en unik identifierare i en nyckel-värde-tabell (Map) som används för att slå upp ett associerat värde; varje nyckel förekommer högst en gång i en Map och kan vara av vilken typ som helst med korrekt equals och hashCode.}
\newcommand{\ExplainKeyvalueTable}{En \textbf{nyckel-värde-tabell} (Map) är en oordnad samling av par där varje unik nyckel mappas till ett värde och som stöder effektiv sökning, insättning och borttagning baserat på nyckeln; i Scala finns den som HashMap, TreeMap m.fl.}
\newcommand{\ExplainLazyInitialization}{\textbf{Lat initialisering} med nyckelordet 'lazy' innebär att en vals värde beräknas och allokeras först vid första åtkomst, inte vid definitionen; detta sparar resurser om värdet aldrig används och löser initialiseringsordningsproblem vid cirkulära beroenden.}
\newcommand{\ExplainLinearSearch}{Att \textbf{linjärsöka} innebär att man traverserar en sekvens element för element tills sökkriteriet är uppfyllt eller alla element undersökts; tidskomplexiteten är O(n) och algoritmen kräver ingen förutgående sortering av data.}
\newcommand{\ExplainLinearSearchAlgorithm}{\textbf{Linjärsökning} är en sökalgoritm som itererar sekventiellt genom en samling och jämför varje element med sökvärdet tills en matchning hittas eller samlingen är slut; algoritmen är enkel men har O(n) tidskomplexitet i värsta fall.}
\newcommand{\ExplainLiteral}{En \textbf{litteral} är en direkt notation för ett specifikt värde i källkoden, t.ex. '42' för ett heltal, '3.14' för ett Double, '"hej"' för en sträng, 'true' för ett booleskt värde, eller ''a'' för ett tecken.}
\newcommand{\ExplainMap}{\textbf{Map}-operationen applicerar en given funktion på varje element i en samling och returnerar en ny samling av samma längd med de transformerade elementen; originalsamlingen förändras inte, vilket gör map till en grundläggande operation i funktionell programmering.}
\newcommand{\ExplainMapping}{En \textbf{mappning} är ett nyckel-värde-par i en Map där nyckeln unikt identifierar värdet; i Scala skrivs en mappning med syntaxen 'nyckel -> värde' och flera mappningar kan kombineras för att konstruera en Map.}
\newcommand{\ExplainMatrix}{En \textbf{matris} är en tvådimensionell datastruktur med m rader och n kolumner indexerad med rad- och kolumnindex; i Scala representeras den ofta som en Vector av Vektorer, och används inom t.ex. bildbehandling och linjär algebra.}
\newcommand{\ExplainMember}{En \textbf{medlem} är en namngiven del av ett objekt eller en klass, t.ex. en metod, ett attribut eller ett typalias; synliga (icke-privata) medlemmar nås utifrån med punktnotation, och kan vara konkreta eller abstrakta.}
\newcommand{\ExplainMemoryComplexity}{\textbf{Minneskomplexitet} beskriver hur mycket minne ett program eller en algoritm behöver som funktion av indatastorleken n, uttryckt med O-notation; O(1) innebär konstant minne, O(n) linjärt och O(n²) kvadratiskt minnesbehov.}
\newcommand{\ExplainMethod}{En \textbf{metod} är en funktion som är definierad som del av ett objekt eller en klass och har tillgång till objektets egna attribut och andra medlemmar; metoder anropas med punktnotation på en objektreferens.}
\newcommand{\ExplainMixin}{\textbf{Inmixning} innebär att man utökar en klass eller ett objekt med egenskaperna hos en eller flera traits; detta möjliggör återanvändning av kodegenskaper och stöder ett slags multipelt arv av beteenden.}
\newcommand{\ExplainModule}{En \textbf{modul} är en avgränsad kodenhet som samlar relaterade abstraktioner under ett gemensamt namn och exponerar ett väldefinierat gränssnitt; i Scala realiseras moduler som singelobjekt, paket eller filer och stöder separation of concerns.}
\newcommand{\ExplainNameShadowing}{\textbf{Namnskuggning} uppstår när ett lokalt namn i ett inre block har samma identifierare som ett namn i ett omgivande block; inom det inre blocket syftar identifieraren på den lokala definitionen och det yttre namnet är otillgängligt utan kvalificering.}
\newcommand{\ExplainNamedArguments}{\textbf{Namngivna argument} låter anroparen ange parameternamnet explicit vid anropet, t.ex. 'f(y = 2, x = 1)'; detta gör att argument kan ges i valfri ordning, förbättrar läsbarheten och fungerar bra i kombination med defaultargument.}
\newcommand{\ExplainNamespace}{En \textbf{namnrymd} är en omgivning där alla namn är unika och kan identifieras entydigt; paket och objekt i Scala skapar egna namnrymder som förhindrar namnkollisioner och gör det möjligt att ha identiska namn i olika moduler.}
\newcommand{\ExplainNew}{Nyckelordet '\textbf{new}' används för att direkt instansiera en klass och anropar dess konstruktor; för case-klasser och objekt med apply-fabriksmetoder behöver 'new' vanligtvis inte skrivas, men krävs vid instansiering av vanliga klasser och anonyma klasser.}
\newcommand{\ExplainNull}{\textbf{Null} är ett specialvärde av referenstyp som indikerar att en referens inte pekar på någon instans; att referera null ger NullPointerException och null rekommenderas inte i Scala — använd istället Option[T] för att representera ett möjligen saknat värde.}
\newcommand{\ExplainObject}{Ett \textbf{objekt} i Scala är antingen ett singelobjekt (definierat med 'object') som samlar variabler och funktioner i en modul och existerar i exakt en upplaga, eller en instans av en klass med eget tillståndsminne skapad med 'new'.}
\newcommand{\ExplainOrdering}{\textbf{Ordning} definierar en total ordning på element av en viss typ via en jämförelsefunktion som returnerar negativt, noll eller positivt tal; Scalas Ordering-typklass används av sorteringsalgoritmer och sorterade samlingar som SortedSet.}
\newcommand{\ExplainOverloading}{\textbf{Överlagring} innebär att en klass eller ett objekt har flera metoder med samma namn men med olika parameterlistor (antal eller typer); kompilatorn väljer rätt variant vid kompilering baserat på argumentens typer, vilket kallas statisk bindning.}
\newcommand{\ExplainOverriddenMember}{En \textbf{överskuggad medlem} är en metod eller val i en subtyp som med nyckelordet 'override' ersätter en namngiven medlem i supertypen; vid dynamisk bindning är det subtypens implementation som anropas, även när referensen är av supertyp.}
\newcommand{\ExplainPackage}{Ett \textbf{paket} i Scala är en namngiven namnrymd som grupperar relaterade klasser och objekt; paket deklareras med 'package' och deras kompilerade bytekod lagras i en katalogstruktur som speglar paketnamnet, t.ex. 'se.lth.cs.prog'.}
\newcommand{\ExplainParameterList}{En \textbf{parameterlista} omges av parenteser och innehåller en kommaseparerad lista av parametrar med namn och typ; en funktion kan ha flera parameterlistor (currying), t.ex. för implicit-parametrar eller för att möjliggöra partiell applikation.}
\newcommand{\ExplainPersistence}{\textbf{Persistens} är egenskapen hos data att överleva bortom en enskild programkörning; data lagras persistent i t.ex. filer, databaser eller andra externa lagringsmedier och kan läsas tillbaka i en senare programkörning.}
\newcommand{\ExplainPolymorphism}{\textbf{Polymorfism} innebär att ett värde kan ha många möjliga typer; i Scala uppnås detta via subtyppolymorfism (arv och överskuggning), parameterpolymorfism (generics och typparametrar) och ad hoc-polymorfism (typklasser och implicita parametrar).}
\newcommand{\ExplainPredicate}{Ett \textbf{predikat} är en funktion returnerar ett booleskt värde (true eller false); predikat används som filterfunktioner i samlingsoperationer som 'filter', 'exists', 'forall' och 'takeWhile'.}
\newcommand{\ExplainPrivate}{En \textbf{privat} medlem, deklarerad med 'private', är synlig endast inuti den klass eller det objekt den tillhör; detta döljer implementationsdetaljer från användare av klassen, skyddar mot godtycklig förändring och minskar kopplingen mellan klasser.}
\newcommand{\ExplainProcedure}{En \textbf{procedur} är en funktion vars syfte är att utföra en sidoeffekt, t.ex. skriva ut text, ändra ett attribut eller kommunicera med omvärlden; returvärdet är av typen Unit som representerar ett tomt värde utan information.}
\newcommand{\ExplainProgramArgument}{\textbf{Programargument} är strängar som skickas till programmet från kommandoraden när det startas; de är tillgängliga via parametern 'args: Array[String]' i main-metoden och används för att anpassa programmets beteende utan att ändra källkoden.}
\newcommand{\ExplainProtectedMember}{En \textbf{skyddad medlem}, deklarerad med 'protected', är synlig inuti den klass den tillhör och i alla direkta och indirekta subklasser, men inte utifrån i övrigt; skyddade medlemmar kan överskuggas i subklasser.}
\newcommand{\ExplainPureFunction}{En \textbf{äkta funktion} (pure function) ger alltid exakt samma returvärde för samma argument och har inga sidoeffekter som påverkar tillståndet utanför funktionen; äkta funktioner är lätta att testa, resonera om, memoize och parallellisera.}
\newcommand{\ExplainRandomSeed}{Ett \textbf{slumptalsfrö} är ett startvärde för en pseudoslumptalsgenerator; med samma frö produceras alltid exakt samma sekvens av pseudoslumptal, vilket möjliggör reproducerbara experiment och deterministiska tester av kod som använder slumpmässighet.}
\newcommand{\ExplainRecursiveFunction}{En \textbf{rekursiv funktion} löser ett problem genom att anropa sig själv med ett enklare eller mindre delfall tills ett basfall nås som kan lösas direkt; svansrekursion (tail recursion) optimeras av Scala-kompilatorn till en loop och undviker stackspill.}
\newcommand{\ExplainReferenceEquality}{\textbf{Referenslikhet} jämför om två referenser pekar på exakt samma objekt i minnet och implementeras i Scala med 'eq'; två separata objekt med identiskt innehåll är inte referenslika, till skillnad från innehållslikhet som jämförs med '=='.}
\newcommand{\ExplainReferenceType}{En \textbf{referenstyp} är en typ vars instanser allokeras på heapen och nås via en referens (pekare); alla klasser i Scala som ärver från AnyRef är referenstyper, t.ex. String, List och egna klasser, till skillnad från värdetyper som Int och Double.}
\newcommand{\ExplainRowVector}{En \textbf{radvektor} är en matris med dimensionen 1×m som innehåller m element i en enda horisontell rad; den används inom linjär algebra och numerisk beräkning och är ett specialfall av en matris.}
\newcommand{\ExplainRuntimeError}{Ett \textbf{exekveringsfel}, även kallat körtidsfel, inträffar medan programmet körs, t.ex. ArithmeticException vid division med noll eller IndexOutOfBoundsException vid felaktigt index; det leder vanligtvis till att programmet avslutas med ett felmeddelande och en stack trace.}
\newcommand{\ExplainRuntimeType}{\textbf{Körtidstyp}en är den faktiska typen hos ett objekt vid exekvering och kan vara mer specifik än den statiska typen som kompilatorn känner till; körtidstypen avgör vilken överskuggad metod som anropas vid dynamisk bindning.}
\newcommand{\ExplainScript}{Ett \textbf{skript} i Scala är en ensam kodfil vars toppnivåkod exekveras direkt utan att en explicit main-metod eller ett omgivande objekt behöver definieras; detta förenklar snabba program, automatisering och interaktiva sessioner.}
\newcommand{\ExplainSealedType}{En \textbf{förseglad typ}, deklarerad med 'sealed', tillåter bara subtyper definierade i samma kodfil; detta gör uppsättningen subtyper stängd och möjliggör uttömmande kontroll vid mönstermatchning, och kompilatorn varnar om ett fall saknas.}
\newcommand{\ExplainSearch}{\textbf{Sökning} är ett algoritmiskt mönster som letar upp ett eller flera element i en samling som uppfyller ett sökkriterium; effektiviteten beror på om data är sorterat och vilken datastruktur som används, t.ex. O(n) för linjärsökning och O(log n) för binärsökning.}
\newcommand{\ExplainSequencecollection}{En sekvenssamling är en ordnad samling med noll eller fler element av samma typ där varje elements position (index) är väldefinierad; element kan nås via index, och sekvenser stödjer operationer som map, filter, foldLeft och zip.}
\newcommand{\ExplainSequenceAlgorithm}{En \textbf{sekvensalgoritm} löser ett problem genom att utnyttja att data är organiserad som en ordnad sekvens; typiska mönster är sökning, sortering, filtrering, mappning, summering och reducering, ofta implementerade med samlingsoperationer eller loopar.}
\newcommand{\ExplainSequenceCollection}{En \textbf{sekvenssamling} är en datastruktur där elementen lagras och nås i en bestämd ordning via heltalsindex; exempel i Scala är Vector (oföränderlig, effektiv), List (oföränderlig, länkad), Array (föränderlig, fast storlek) och ArrayBuffer (föränderlig, dynamisk storlek).}
\newcommand{\ExplainSerialize}{Att \textbf{serialisera} innebär att man kodar ett objekt till en sekvens av tecken eller bytes i ett väldefinierat format (t.ex. JSON, XML eller ett binärt format) som kan lagras på disk, skickas över ett nätverk och sedan deserialiseras tillbaka till ett objekt.}
\newcommand{\ExplainSet}{En \textbf{mängd} (Set) är en oordnad samling av unika element utan duplicat; Scala erbjuder både oföränderliga och föränderliga implementationer med operationer som union, snitt, differens och effektiv membership-test.}
\newcommand{\ExplainSetter}{En \textbf{setter} är en metod som ger indirekt skrivåtkomst till ett privat attribut; i Scala kan en setter definieras med syntaxen 'def attributnamn_=(värde: Typ)' vilket möjliggör tilldelning med vanlig '='-syntax från utsidan av klassen.}
\newcommand{\ExplainSingletonObject}{Ett \textbf{singelobjekt} definieras med nyckelordet 'object' och existerar i exakt en upplaga per program; det kan ha tillståndsvariabler och metoder, initialiseras vid första åtkomst, och används som modul, kompanjonsobjekt eller programingångspunkt.}
\newcommand{\ExplainSorting}{\textbf{Sortering} är ett algoritmiskt mönster som ordnar elementen i en samling i stigande eller fallande ordning enligt en given ordning; effektiva algoritmer som mergesort och quicksort har O(n log n) tidskomplexitet, och Scala erbjuder inbyggd sortering via 'sorted' och 'sortBy'.}
\newcommand{\ExplainStackTrace}{En \textbf{stack trace} är en lista som skrivs ut när ett körtidsfel inträffar och som beskriver den fullständiga anropskedjan från felpunkten upp till programmets startpunkt, med information om klassnamn, metodnamn, filnamn och radnummer.}
\newcommand{\ExplainStatement}{En \textbf{sats} är en kodkonstruktion som utför en åtgärd, t.ex. en tilldelning, ett metodanrop eller en kontrollstruktur; i Scala är satser även uttryck och har ett värde (ofta Unit), och flera satser på samma rad separeras med semikolon.}
\newcommand{\ExplainString}{En \textbf{sträng}  är en oföränderlig sekvens av Unicode-tecken; i Scala är String ett alias för java.lang.String och erbjuder ett rikt urval av metoder för sökning, ersättning, delsträngar och transformation.}
\newcommand{\ExplainStringInterpolator}{Att 'interpolera' innebär att skjuta in mellanliggande ord i en text. Stränginterpolation är en Scala-funktion som tillåter inbäddning av variabelreferenser och uttryck direkt i strängliteraler med hjälp av s-interpolatorn, som skrives med ett inledande s-tecken.}
\newcommand{\ExplainStructuralEquality}{\textbf{Innehållslikhet} (strukturell likhet) innebär att två instanser anses lika om deras fältvärden är lika; i Scala implementeras detta med 'equals' och '==', och genereras automatiskt av case-klasser; vanliga klasser ärver referenslikhet från AnyRef som standard.}
\newcommand{\ExplainSubtype}{En \textbf{subtyp} är en typ som är mer specifik än sin supertyp och kan användas överallt där supertypen förväntas (Liskovs substitutionsprincip); en subtyp ärver supertypens gränssnitt och kan lägga till nya medlemmar eller överskugga befintliga.}
\newcommand{\ExplainSupertype}{En \textbf{supertyp} är en typ som är mer generell och vars gränssnitt ärvs av subtyper; en variabel av supertyp kan referera till instanser av vilken subtyp som helst, vilket möjliggör polymorfism och utbytbarhet av implementationer.}
\newcommand{\ExplainTimeComplexity}{\textbf{Tidskomplexitet} beskriver hur exekveringstiden för en algoritm växer som funktion av indatastorleken n, uttryckt med O-notation; t.ex. är O(1) konstant, O(n) linjär, O(n log n) typisk för effektiv sortering och O(n²) kvadratisk tidskomplexitet.}
\newcommand{\ExplainTrait}{En \textbf{trait} är en abstrakt typ som kan ha både abstrakta och konkreta medlemmar; den kan mixas in i klasser med 'with', kan i Scala 3 ha konstruktorparametrar, och möjliggör ett slags multipelt arv av beteenden utan diamantproblemet.}
\newcommand{\ExplainType}{En \textbf{typ} är en etikett på data som beskriver vilka operationer som är tillåtna och vilka värden som är möjliga; Scalas statiska typsystem kontrollerar typkorrekthet vid kompilering och förhindrar många klasser av fel innan programmet körs.}
\newcommand{\ExplainTypeAlias}{Ett \textbf{typalias}, definierat med 'type Namn = BefintligTyp', skapar ett alternativt namn för en typ och ökar kodens läsbarhet och domänuttryckskraft utan att påverka körbeteendet; typalias kan parameteriseras och kan vara abstrakta i traits.}
\newcommand{\ExplainTypeArgument}{Ett \textbf{typargument} är en konkret typ som anges inom hakparenteser när en generisk klass eller funktion används, t.ex. 'Vector[Int]' där 'Int' är typargumentet; kompilatorn binder typargumentet till typparametern och kontrollerar att det är korrekt använt.}
\newcommand{\ExplainTypeInference}{\textbf{Typhärledning} innebär att Scala-kompilatorn automatiskt bestämmer typen för ett uttryck eller en variabel utifrån sammanhanget, så att explicita typannoteringar ofta kan utelämnas; detta minskar kodmängden utan att försaka typsäkerheten.}
\newcommand{\ExplainUniformAccess}{Principen om \textbf{enhetlig access} innebär att en egenskap kan implementeras antingen som ett attribut (val/var) eller en metod (def) utan att anropande kod behöver ändras; detta ger frihet att byta implementation utan att påverka klassens gränssnitt.}
\newcommand{\ExplainValueType}{En \textbf{värdetyp} är en typ vars instanser lagras direkt på stacken eller inline i ett objekt och kopieras vid tilldelning; i Scala är Int, Double, Boolean, Char m.fl. värdetyper med supertypen AnyVal och boxas automatiskt till objekt vid behov.}
\newcommand{\ExplainVariable}{En \textbf{variabel} används för att ge datavärden ett namn. I de delar av koden där variabeln är synlig kan värdet refereras via variabelns namn. En variable deklareras med val eller var och ges då ett initialvärde. En val-variabel är oföränderlig och kan ej ändras efter initialisering medan en var-variabel kan tilldelas nya värden efter initialisering genom tilldelning.}
\newcommand{\ExplainWhileStatement}{En \textbf{while-sats} upprepar sin kropp så länge ett booleskt villkor är sant; antalet iterationer behöver inte vara känt i förväg, vilket gör den lämpad för algoritmer som körs tills ett villkor är uppfyllt, t.ex. inläsning tills filen är slut.}
\newcommand{\ExplainYield}{Nyckelordet '\textbf{yield}' används i ett for-uttryck för att samla upp de producerade värdena i en ny samling; typen på den nya samlingen är densamma som den iterade samlingens typ, och for-yield är syntaktisk socker för flatMap/map-kombinationer.}